\scanpage{95}
% loc 1–7
Il faudrait prouver que ça implique qu'il y a un $x\in k.1\oplus ku\oplus kv$, partout non nul, qui est orth. (pour $\operatorname{Tr}(xy)$) à $1,u,v$ -- il faut donc calculer le \deleted{déterminant} discriminant de $\operatorname{Tr}$ par rapport à la base $1,u,v$, on trouve
% loc 8–11
\[
\begin{aligned}
&\text{discriminant}\\[-2pt]
&\text{de }\operatorname{Tr}\text{ pour}\\[-2pt]
&\text{la base }1,u,v\\[-2pt]
&\text{de }k1\oplus ku\oplus kv
\end{aligned}
=-2\,Q(t;b,c,b',c')
\]
% loc 12–14, bracketed addition
[ et le discriminant \uline{réduit} de la forme \uline{quadratique} $\det$ est
\[ -Q(t;b,c,b',c')\qquad\biggr] \]
% loc 12–15, left margin
% slanted
\margin{? Vérifier si c'est ça \uuline{Oui}}
% loc 15–17
Donc le fait que c'est nul \deleted{signifie}\inserted{implique} que \deleted{\unsure{les}} $\exists\ \unsure{\theta|\Theta}$ orthogonal à $1,u,v$, partout non nul et \uuline{isotrope}

% loc 18–19
Donc \deleted{que}
\[ \det\theta=\operatorname{Tr}\theta=0 \]
% loc 20–22
Or \deleted{l'une des x} les droites isotropes dans l'hyperplan $\operatorname{Tr}\theta=0$ -- correspondent aux Borels, \deleted{comme les} on gagne !
% archive sheet number: 94
